\section{Mechanistic Interpretability 基础}

\subsection{什么是 Mech Interp？}

神经网络是``生长的而非编程的''——架构是脚手架，损失函数是指引光源，梯度下降产生了一个我们不知道如何直接编程的产物。\textit{``我们不是制造它们，而是培育它们……这几乎是一个生物实体或有机体，我们在研究它。''}

这里有两个相互交织的问题：一个深刻的科学问题（这些系统内部到底发生了什么？），和一个关键的安全问题（如何确保它们可信？）。

\begin{importantbox}{不是归因，是逆向工程}
Saliency Map（例如``图像的哪个部分让模型认为这是狗''）\textbf{不是} Mechanistic Interpretability。Mech Interp 追求的是\textbf{算法和机制}——将神经网络权重逆向工程为可理解的算法。

类比：神经网络权重 = 二进制计算机程序。目标是将这个``编译后的程序''反编译成可读的算法。激活值 = 内存；权重 = 指令；两者都需要理解。

核心态度：\textit{``Gradient descent 比你聪明''} ——自下而上发现，而非自上而下假设。不要预设模型内部有什么，去发现它。
\end{importantbox}

\subsection{Universality：跨网络甚至跨物种的一致性}

不同架构的视觉模型都会形成相同的特征：Gabor 滤波器、曲线检测器、高低频检测器。更令人震惊的是，这些特征也存在于\textbf{生物}神经网络中：

\begin{knowledgebox}{跨物种的特征一致性}
\begin{itemize}
    \item \textbf{曲线检测器}：先在人工神经网络中发现，后在猴子大脑中确认
    \item \textbf{高低频检测器}：先在人工神经网络中发现，后在小鼠大脑中确认
    \item \textbf{Donald Trump 神经元}：每个视觉模型都有一个专门的 Trump 神经元——同时响应他的面部图像\textbf{和}``Trump''这个词。这是一个抽象概念，不仅仅是模式匹配
\end{itemize}

\textit{``梯度下降在某种意义上找到了正确的方式来切分事物……许多系统都收敛到相同的抽象。''} 这暗示存在某种``自然的抽象边界''，无论计算基质是硅还是碳，都会被发现。
\end{knowledgebox}

\subsection{本章小结}

Mechanistic Interpretability 追求理解神经网络的``算法''，而非仅仅归因。跨网络和跨物种的特征一致性暗示梯度下降在发现某种``自然的抽象切分方式''。

